The thesis has six chapters.

\begin{itemize}
    \item Chapter 1 presents the motivations for exploring augmentation techniques, the problem statement, and the overall structure of the thesis.
    \item Chapter 2 serves as a source of relevant background information for the thesis, including an overview of related works, a detailed description of the marine wildlife dataset, and an introduction to important concepts in Generative Adversarial Network (GAN).
    \item Chapter 3 of the thesis covers the fundamental concepts of basic augmentation techniques, providing a detailed description of the basic augmentations and a systematic approach to tuning their hyperparameters.
    \item Chapter 4 outlines the experimental approach, including the evaluation metrics and models employed, as well as the training pipelines. The training pipelines are divided into two categories: the experiment pipeline for basic augmentations and the GAN-generated pipeline, where the specific architectures of GANs will be discussed.
    \item Chapter 5 presents the outcomes of the experiments and subsequent analysis.
    \item Chapter 6 summarized the findings from the experiments, outlines the limitations of the thesis, and proposes potential avenues for future research.
\end{itemize}










